\subsection{Exponential and Trigonometic Functions}

A power series with radius of convergence $R \in (0, \infty]$ gives a function $f:(-R,R) \rightarrow \mathbb{R}$. Many important functions can be defined in this way.

\begin{definition}
    For $x \in \mathbb{R}$, define
    \[\exp(x) := \sum_{n=0}^{\infty} \frac{x^n}{n!}.\]
\end{definition}

\begin{prop}
    The power series $\sum_{n=0}^{\infty} \frac{x^n}{n!}$ has infinite radius of convergence, thus the exponential function has domain equal to $\mathbb{R}$.
\end{prop}

\begin{prop}
    \[\exp(x + y) = \exp(x)\exp(y).\]
\end{prop}
\begin{proof}
    Non-examinable.
\end{proof}

\begin{corollary}
    For any $x \in \mathbb{R}$,
    \begin{itemize}
        \item $\exp(x) > 0$
        \item $\exp( - x) = \frac{1}{\exp(x)}$.
    \end{itemize}
\end{corollary}
\begin{proof}
    \begin{itemize}
      \item We know $\exp(0) = 1$. Then
        \begin{align*}
          \exp(x - x) &= 1 \\
              \exp(x) \exp(-x) &= 1
              \exp(-x) &= \frac{1}{\exp(x)}
        .\end{align*}
    \item If $x \geq 0$, $\exp(x) \geq 1 > 0$. Then
      \[\exp(-x) = \frac{1}{\exp(x)} > 0.\]
    \end{itemize}
\end{proof}


\begin{definition}
    For all $x \in \mathbb{R}$,
    \[\cos(x) = \sum_{n=0}^{\infty} \frac{(-1)^n}{(2n)!} x ^ {2n}.\]
    \[\sin(x) =\sum_{n=0}^{\infty} \frac{(-1)^n}{(2n + 1)!} x ^{2n + 1}.\]
\end{definition}

\begin{remark}
    We cannot find the radius of convergence using the ratio test, as many terms are equal to $0$. However, we can compare these series with the exponential function and show they are bounded by the comparison test.
\end{remark}
